\scanpage{7}
donc la relation précédente équivaut à
\[
p^{i+1}\Theta_{i+1}=0
\]
On a donc prouvé les relations $p^i\Theta_i=0$. Pour la deuxième relation, on écrit
\[
(\lambda_\nu+\Theta_\nu)^p=0\qquad\text{i.e.}
\]
\[
\lambda_\nu^p+\sum_{i=1}^p C_i^p\Theta_\nu^i\lambda_\nu^{p-i}=0
\]
Multiplions par $p^\nu$ et utilisons que $p^\nu\Theta_\nu\in I$ donc $p^\nu\Theta_\nu^\alpha\in I$ si $\alpha\geq1$, on trouve
\[
p^\nu\lambda_\nu^p\in I
\]
Or on a
\[
v_p(p^\nu\lambda_\nu^p)=\nu+p\,v_p\left(\underbrace{\lambda_\nu}_{\inserted{=\frac{p^{p^\nu}}{p^\nu!}}}\right)=\nu+p\left(p^\nu-\frac{p^\nu-1}{p-1}\right)=c(\nu)
\]
Pour prouver que la condition est nécessaire, il suffit de noter que l'ensemble $I'$ formé des éléments satisfaisant les cond. \deleted{\unsure{de}} ci-dessus données (qui satisfait $I'\subset I$, $\Theta_0\in I'$) est bien un \uline{idéal}, i.e. \deleted{est stable} (comme c'est un sous-$W$-module) est stable par multiplication par $\Theta_0,\ldots,\Theta_\nu$. La vérification est facile et laissée au lecteur.

\margin{!}

\uline{Corollaire} Soit, pour tt $i$, $E_i$ un syst. de représentants de $W/p^iW$. Alors un élément de $W^{(\nu)}$ s'écrit de façon unique sous la forme
\[
\sum w_{a_1,\ldots,a_\nu}\dot{\Theta}_1^{a_1}\cdots\dot{\Theta}_\nu^{a_\nu}
\]
avec
\[
\left\{\begin{aligned}
&0\leq a_i\leq p-1\qquad(1\leq i\leq\nu)\\
&w_{a_1,\ldots,a_\nu}\in E_i\quad\text{si }a_i\text{ est la première coordonnée non nulle de }(a_1,\ldots,a_\nu).
\end{aligned}\right.
\]
\mspage{6}
